\newcommand{\speechpart}{Part C and closing, slides 16--20}
\documentclass[11pt,a4paper]{article}
\usepackage[utf8]{inputenc}
\usepackage[T1]{fontenc}
\usepackage{lmodern}
\usepackage[margin=2.2cm]{geometry}
\usepackage{xcolor}
\usepackage{titlesec}
\usepackage{fancyhdr}
\usepackage[most]{tcolorbox}
\usepackage{microtype}
\usepackage{needspace}
\usepackage[hidelinks]{hyperref}
\definecolor{unicablue}{RGB}{45,58,104}
\definecolor{unicared}{RGB}{153,38,45}
\definecolor{tint}{RGB}{228,231,240}
\setlength{\parindent}{0pt}
\setlength{\parskip}{6pt}
\linespread{1.15}
\titleformat{\section}{\Large\bfseries\color{unicablue}}{}{0pt}{}
\titlespacing*{\section}{0pt}{12pt}{2pt}
\newcommand{\slide}[2]{\Needspace{10\baselineskip}\section{#1}\textcolor{unicared}{\small\bfseries #2}\par}
\newcommand{\cue}[1]{{\small\itshape\color{unicared}[#1]}}
\newcommand{\pause}{\textcolor{unicared}{\,/\!/\,}}
\newtcolorbox{say}{breakable,colback=white,colframe=unicablue,boxrule=0pt,leftrule=3pt,arc=0pt,
  left=10pt,right=6pt,top=4pt,bottom=4pt,fontupper=\large,before upper={\setlength{\parskip}{8pt}}}
\pagestyle{fancy}\fancyhf{}
\fancyhead[L]{\small\color{unicablue}Measuring Agglomeration --- \speechpart}
\fancyhead[R]{\small\color{unicablue}4 November 2026}
\fancyfoot[C]{\small\thepage}

\begin{document}
{\color{unicablue}\Huge\bfseries Speech: slides 16--20\par}
\vspace{2pt}
{\Large Part C --- the Berlin Wall (Ahlfeldt, Redding, Sturm and Wolf 2015) --- and closing\par}
\vspace{4pt}
{\small About 8 minutes 30 seconds. \cue{red brackets} = what to point at or do; \pause\ = a short pause. Numbers from the updated report (ARSW2). If Part C is drawn, open with slides 1--2 (Part A speech, slides 1--7), then say: ``Let me go to Part C.''}

\slide{Slide 16 --- OLS finds three times the true lambda, even where there is none}{about 1 min 30 s}
\begin{say}
\cue{title} Part C: the Berlin Wall. Here lambda is a spillover between city blocks: productivity depends on the density of employment around each block, Upsilon. Point (a). We ran ARSW0 as it is: it simulates a city where the true lambda is \textbf{0.07}, and regresses log productivity on log surrounding density.

\cue{table} With random fundamentals, OLS gives \textbf{0.206} on average over 20 draws; when fundamentals favour the centre, \textbf{0.218}. \pause\ Both are about three times the true value. \cue{figure} In both panels the red OLS line is much steeper than the dashed true slope.

\cue{formula} Why? Log productivity is log A-bar plus lambda log Upsilon. A-bar is unobserved, so it sits in the \textbf{error term}, and it is correlated with density: the covariance is positive, so lambda-hat is biased upward. \cue{left card} In case (a) fundamentals are random in space, so the bias is pure \textbf{reverse causality}: a block with a good fundamental attracts workers, so the density around it rises --- density responds to productivity instead of causing it. In case (b) there is also an \textbf{omitted variable}, the centre is denser and better, but it adds only 0.012.

\cue{right card} Policy. Doubling density gives \textbf{4.97 percent} with the true value, but OLS promises \textbf{15.35}. \cue{last row of the table} And in a city with \emph{no} spillover, OLS still reports 0.196: the regression cannot tell the two cities apart. A government would expect 14.6 percent where the true gain is zero, and subsidise concentration where there is no market failure to correct.
\end{say}

\slide{Slide 17 --- The division of the city recovers the true lambda}{about 1 min 30 s}
\begin{say}
\cue{title} Point (b): act 5 of ARSW2, the professor's updated script, run as it is. The Wall takes away the employment around the blocks near the old centre, without changing what those blocks are.

How does the estimator work? For each candidate lambda, it inverts the model \textbf{twice}, before and after the Wall, and recovers the fundamentals at both dates. Then it takes the change in log A-bar in the West blocks and regresses it on \textbf{distance to the old centre}. \pause\ The estimate is the lambda at which that slope is \textbf{zero}.

\cue{figure and navy box} With fundamentals unchanged, the blue line crosses zero exactly at \textbf{0.0700}, the true value. \cue{table} With random drift in fundamentals, five draws give 0.053 to 0.087: mean \textbf{0.071}. Five draws are a suggestive check, not a proof, but they scatter around the true value.

\cue{first card} The identifying assumption: the Wall changed what blocks could \emph{reach}, not what they \emph{are}: the change in the fundamental A-bar is uncorrelated with distance to the old centre.

\cue{second card} The moment at lambda zero is \textbf{plus 0.0204}. It is a slope: positive means fundamentals fall near the centre. The central blocks really lost productivity, because they lost density; with lambda zero the model cannot blame density, so it reads the loss as a fall in the fundamentals --- for no reason. So lambda zero is rejected.

\cue{third card} Why does policy need this and not OLS? OLS gives about 0.20 whether the true lambda is 0.07 or zero. The Wall holds fundamentals fixed and isolates the true spillover.
\end{say}

\slide{Slide 18 --- A wrong lambda overstates the welfare gain of the metro by 56\%}{about 2 min}
\begin{say}
\cue{code} Point (c). Act 6 builds a metro line along the central row: trains at 3 minutes per kilometre, 2 minutes of waiting, and every commuting pair takes the faster of walking and the metro. Our only change to the code is \texttt{lambda\_policy}: we ran the act with zero, 0.07 and 0.18. For each value the model recovers the fundamentals from the \emph{same} observed city, so before the line the three cities are identical; only the prediction after the line differs. \cue{footnote} At lambda zero only the production spillover is removed; the residential externality stays on.

\cue{table} In the table, blue is an increase and red a decrease.

\cue{point (i)} Who gains? About a third of commuting pairs get faster, \textbf{6,476 of 20,449}, saving a median of \textbf{4.4 minutes}: commuters along the corridor, and owners of floor space near the line. In spatial equilibrium the saving is \textbf{capitalised} in floor prices: they rise on the line by 13 to 32 percent, and fall away from it by 2 to 3 percent.

\cue{point (ii), jobs row} Central jobs. \pause\ With lambda zero the metro only makes commuting cheaper: workers can live farther away, nothing anchors firms in the centre, and central jobs \textbf{fall} by 3.85 percent. With lambda 0.07, more workers reach the centre, density rises, and density raises productivity there --- the same feedback loop as in Part A: central jobs \textbf{rise} by 2.49 percent. The sign flips. At 0.18 the loop is so strong that central jobs jump by 52 percent and firms crowd residents out of the line, minus 35.

\cue{point (iii), welfare row, look at the Professor} The wrong lambda. A government using the OLS value would predict a welfare gain of \textbf{2.60 percent} instead of \textbf{1.67}: 1.56 times, a \textbf{56 percent} overstatement, 0.93 percentage points. And it is conservative: the updated OLS estimate is actually 0.206.
\end{say}

\slide{Slide 19 --- 0.07 is a block-level number}{about 1 min 15 s}
\begin{say}
\cue{navy banner} Point (d), optional. Would we use 0.07 for \texttt{par.lambda} in ARM2? At most as an upper bound. 0.07 is an \textbf{intra-city} estimate: with delta 0.36 per minute, the spillover \textbf{halves every 1.9 minutes}, about 160 metres on foot, so it is negligible beyond a few hundred metres. ARM2 links cities hundreds of kilometres apart.

\cue{``For''} For: it is a standard benchmark and, without regional estimates, an upper-bound scenario for policy; some channels, like thick labour markets and shared suppliers, reach beyond the city.

\cue{``Against''} Against, four arguments. \textbf{Decay}: lambda is inseparable from delta. \textbf{Mechanisms}: blocks interact face to face, regions through freight and trade. \textbf{The model}: in ARM2 a lambda of one over sigma minus one, 0.25 to 0.33, is a Helpman variety gain, as in Part B(d), a different object. \textbf{Evidence}: regional estimates are lower, 0.02 to 0.05, Rosenthal and Strange 2004, Combes and Gobillon 2015.

\cue{``What we would need''} To decide we would need a decay function for inter-regional distances, lambda estimated on regional data, and measured trade costs, to separate spillovers from the market access ARM2 already captures.
\end{say}

\slide{Slide 20 --- Takeaways}{about 1 min}
\begin{say}
To sum up.

\cue{01} \textbf{First}, Part A: a cross-section identifies A, not lambda. The model recovers productivity exactly, but not how much comes from density. And lambda matters: the same shock moves 26.9 percent more people without spillovers, 45.3 percent with lambda 0.25.

\cue{02} \textbf{Second}, Part B: shocks to market access move people. Exposure to the East ranks the losers of the division; mobility and agglomeration change the size of the effect, not the ranking. And Redding and Sturm's implicit lambda of one third is Helpman's one over sigma minus one.

\cue{03} \textbf{Third}, Part C: \pause\ lambda needs a natural experiment. OLS gives about 0.21 for a true 0.07, and 0.20 even with no spillover; the Wall recovers 0.07; and with the OLS value the welfare gain of a metro is overstated by 56 percent.

\cue{look at the Professor} Thank you for your attention. We are happy to take your questions.
\end{say}

\vspace{6pt}
\begin{tcolorbox}[colback=tint,colframe=tint,arc=0pt]
\small\textbf{Numbers to remember.} OLS 0.206 (0.177--0.247), 0.218, 0.196 at $\lambda=0$; $+0.012$ from the omitted variable \ $\cdot$\ $4.97\%$ vs $15.35\%$ ($14.6\%$ at $\lambda=0$) \ $\cdot$\ $\hat\lambda = 0.0700$; drift 0.053--0.087, mean 0.071 \ $\cdot$\ moment $+0.0204$ \ $\cdot$\ 6,476 of 20,449 pairs, 4.4 min \ $\cdot$\ welfare $+1.37$ / $+1.67$ / $+2.60\%$; jobs $-3.85$ / $+2.49$ / $+52.08\%$ \ $\cdot$\ $2.60/1.67 = 1.56$ ($+0.93$ pp) \ $\cdot$\ half-life 1.9 min $\approx$ 160 m \ $\cdot$\ regional 0.02--0.05.

\textbf{Words to get right.} case (a) is \emph{only} reverse causality \ $\cdot$\ the assumption is about $\bar A$ (fundamental), not $A$ \ $\cdot$\ the moment is a \emph{slope} \ $\cdot$\ a violation (e.g.\ slower reconstruction near the centre) biases $\hat\lambda$ \emph{upward} \ $\cdot$\ \emph{intra-city}, not intercity \ $\cdot$\ welfare = expected utility.
\end{tcolorbox}
\end{document}
